\subsection{Destruction}\label{sec:destruction}
The destruction stage reestablishes control flow by extracting an IPG from an RVSDG
and generating CFGs from individual \lambdaregions. Inter-Procedural Control Flow Recovery
(Inter-PCFR) creates an IPG from \lambdanodes, while
Intra-Procedural Control Flow Recovery (Intra-PCFR) extracts control flow from \gammam and
\thetanodes and generates basic blocks with corresponding operations for primitive nodes.
A \lambdaregion without \gammam and \thetanodes trivially transforms into a linear CFG,
while \lambdaregions with these nodes require construction of branches and/or loops.
This section discusses Inter-PCFR in detail. Detailed discussion
of Intra-PCFR is found in Bahmann~\etal~\cite{Bahmann15}.

\subsubsection{Inter-Procedural Control Flow Recovery}\label{sec:ipcfr}
Inter-PCFR recovers an IPG from an RVSDG. IPG nodes are created for \lambdanodes as well as
arguments of the \omegaregion, while IPG edges are inserted to capture the dependencies between
\lambdanodes. Algorithm~\ref{alg:intrapt} starts by
creating IPG nodes for all arguments of the \omegaregion, \ie, all external functions. It
continues by recursively traversing the region tree, creating IPG nodes for encountered
\lambdanodes and IPG edges for their dependencies. For the region of every \lambdanode, it
invokes Intra-PCFR to create a CFG.

\begin{informalg}{Inter-Procedural Control Flow Recovery}{alg:intrapt}
\begin{enumerate}
	\vspace*{-0.7\baselineskip}

	\item Create IPG nodes for all arguments of the \omegaregion.
	\item Process all nodes of the \omegaregion in topological order as follows:
		\begin{itemize}
			\item \bpoint{\lambdanodes}: Create an IPG node, and mark it exported if the
				\lambdanode's output has a \omegaregion's result as user. For every context variable
				$cv = (\iput{}{}, \argument{}{})$, add an edge from the \lambdanode's IPG node to the
				corresponding IPG node of the producer of $\iput{}{}$. Create a CFG from the \lambdanode's
				subregion and attach it to the IPG node.

			\item \bpoint{\deltanodes}: Create an IPG node, and mark it exported if the
				\deltanode's output has a \omegaregion's result as user. For every context variable
				$cv = (\iput{}{}, \argument{}{})$, add an edge from the \deltanode's IPG node to the
				corresponding IPG node of the producer of $\iput{}{}$. Create the expression from the
				\deltanode's subregion and attach it to the IPG node.

			\item \bpoint{\phinodes}: For every argument of the \phiregion, create an IPG node for
				the corresponding \deltam/\lambdanode and add IPG edges from this node to the
				corresponding IPG nodes of the context variables. Translate the \deltam/\lambdanodes in the
				\phiregion according to the rules above. Mark the IPG node as exported if the corresponding
				\phinode's output has a \omegaregion's result as user.
		\end{itemize}

	\vspace*{-\baselineskip}
\end{enumerate}
\end{informalg}

\subsubsection{Intra-Procedural Control Flow Recovery}
Bahmann~\etal~\cite{Bahmann15} explored two different approaches for CFG generation:
\textit{Structured Control Flow Recovery (SCFR)} and \textit{Predicative Control Flow Recovery
(PCFR)}. SCFR uses the region hierarchy within a \lambdaregion to recover control flow, while PCFR
generates branches for predicate producers and follows the predicate consumers to the eventual
destination. Both schemes reestablish evaluation-equivalent CFGs, but differ in the recoverable
control flow. SCFR recovers only control flow that resembles the structural nodes in
\lambdaregions, \ie, control flow equivalent to \verb;if-then-else;, \verb;switch;, and
\verb;do-while; statements, while PCFR can recover arbitrary complex control flow, \ie, control
flow that is not restricted to RVSDG constructs. PCFR reduces the number of static branches in the
resulting control flow~\cite{Bahmann15}, but might also result in undesirable
control flow for certain architectures, such as graphic processing units~\cite{Reissmann16}.
For the sake of brevity, we omit a discussion of SCFR and PCFR as the algorithms are extensively
described by Bahmann~\etal~\cite{Bahmann15}.



